% GENERATED FILE. Edit canonical lessons, then run npm run generate:books.
% canonical-source-hash: fnv1a64:3d5c3e8c3b721c9d
% canonical-lessons: TA-C218-ottakaccivinki, TA-C218-varikkutirai, TA-C218-mutalai, TA-C218-pattampucci, TA-C218-teni

\chapter{Giraffe, Zebra, Crocodile, Butterfly, Bee}
\label{ch:ta-giraffe-zebra-crocodile-butterfly-bee}
\hlchaptermodality{218}

\begin{chapteropening}
\textbf{By the end of this chapter:} \emph{I can say a giraffe, a zebra, a crocodile, a butterfly and a bee.}

Complete the last lesson of chapter 218: I can say a giraffe, a zebra, a crocodile, a butterfly and a bee.
\end{chapteropening}

\section[\texorpdfstring{\ta{ஒட்டகச்சிவிங்கி} (oṭṭakacciviṅki)}{ஒட்டகச்சிவிங்கி (oṭṭakacciviṅki)}]{\ta{ஒட்டகச்சிவிங்கி} (oṭṭakacciviṅki) — a giraffe}

\label{lesson:TA-C218-ottakaccivinki}



\begin{quote}
Before the new one: say the Tamil for a feast, then the Tamil for a festival.
\end{quote}

\subsection*{You'll want to know: \ta{ஒட்டகச்சிவிங்கி}}
\textbf{\ta{ஒட்டகச்சிவிங்கி}} — \emph{oṭṭakacciviṅki} — \textquotedblleft{}a giraffe\textquotedblright{}.

\begin{cousinweb}
\textbf{\ta{ஒட்டகச்சிவிங்கி}} (\emph{oṭṭakacciviṅki}) is \textquotedblleft{}the camel-leopard\textquotedblright{}: \emph{oṭṭakam}, \textquotedblleft{}camel\textquotedblright{}, from Sanskrit \emph{uṣṭra} (CDIAL 2387), and native \emph{civiṅki}, \textquotedblleft{}hunting leopard\textquotedblright{} (DEDR 2579; Malayalam \emph{civiṅṅi}). Greek named the giraffe the same way, \emph{camelopardalis}.
\end{cousinweb}

\begin{grammarlens}[title={What you've built}]
One more word for this chapter.
\end{grammarlens}

\subsection*{Your turn}
\begin{itemize}
\raggedright
  \item \emph{Say it:} \emph{oṭṭakacciviṅki}
  \item \emph{Say it:} \emph{oṭṭakacciviṅki}, once more
  \item \emph{Say it:} say \emph{oṭṭakacciviṅki}
  \item \emph{Recall:} say the Tamil for a feast, then the Tamil for a festival, then say \emph{oṭṭakacciviṅki} again
\end{itemize}

\begin{tcolorbox}[breakable,colback=teal!4,colframe=teal!35!black,title={Before you move on}]
What does \ta{ஒட்டகச்சிவிங்கி} mean? (\textquotedblleft{}A giraffe\textquotedblright{}.) Say it once more.
\end{tcolorbox}
\section[\texorpdfstring{\ta{வரிக்குதிரை} (varikkutirai)}{வரிக்குதிரை (varikkutirai)}]{\ta{வரிக்குதிரை} (varikkutirai) — a zebra}

\label{lesson:TA-C218-varikkutirai}



\begin{quote}
Before the new one: say the Tamil for a festival, then the Tamil for a giraffe.
\end{quote}

\subsection*{You'll want to know: \ta{வரிக்குதிரை}}
\textbf{\ta{வரிக்குதிரை}} — \emph{varikkutirai} — \textquotedblleft{}a zebra\textquotedblright{}.

\begin{cousinweb}
\textbf{\ta{வரிக்குதிரை}} (\emph{varikkutirai}) is native: \emph{vari}, \textquotedblleft{}stripe, line\textquotedblright{} (DEDR 5263), and \emph{kutirai}, \textquotedblleft{}horse\textquotedblright{} (DEDR 1711): \textquotedblleft{}the striped horse\textquotedblright{}.
\end{cousinweb}

\begin{grammarlens}[title={What you've built}]
One more word for this chapter.
\end{grammarlens}

\subsection*{Your turn}
\begin{itemize}
\raggedright
  \item \emph{Say it:} \emph{varikkutirai}
  \item \emph{Say it:} \emph{varikkutirai}, once more
  \item \emph{Say it:} say \emph{varikkutirai}
  \item \emph{Recall:} say the Tamil for a festival, then the Tamil for a giraffe, then say \emph{varikkutirai} again
\end{itemize}

\begin{tcolorbox}[breakable,colback=teal!4,colframe=teal!35!black,title={Before you move on}]
What does \ta{வரிக்குதிரை} mean? (\textquotedblleft{}A zebra\textquotedblright{}.) Say it once more.
\end{tcolorbox}
\section[\texorpdfstring{\ta{முதலை} (mutalai)}{முதலை (mutalai)}]{\ta{முதலை} (mutalai) — a crocodile}

\label{lesson:TA-C218-mutalai}



\begin{quote}
Before the new one: say the Tamil for a giraffe, then the Tamil for a zebra.
\end{quote}

\subsection*{You'll want to know: \ta{முதலை}}
\textbf{\ta{முதலை}} — \emph{mutalai} — \textquotedblleft{}a crocodile\textquotedblright{}.

\begin{cousinweb}
\textbf{\ta{முதலை}} (\emph{mutalai}) is native Dravidian. DEDR 4952 lists Malayalam \emph{mutala}, Kannada \emph{mosaḷe} and Telugu \emph{mosali}. The entry compares Sanskrit (lexicons) \emph{mācala}, \textquotedblleft{}crocodile\textquotedblright{}.
\end{cousinweb}

\begin{grammarlens}[title={What you've built}]
One more word for this chapter.
\end{grammarlens}

\subsection*{Your turn}
\begin{itemize}
\raggedright
  \item \emph{Say it:} \emph{mutalai}
  \item \emph{Say it:} \emph{mutalai}, once more
  \item \emph{Say it:} say \emph{mutalai}
  \item \emph{Recall:} say the Tamil for a giraffe, then the Tamil for a zebra, then say \emph{mutalai} again
\end{itemize}

\begin{tcolorbox}[breakable,colback=teal!4,colframe=teal!35!black,title={Before you move on}]
What does \ta{முதலை} mean? (\textquotedblleft{}A crocodile\textquotedblright{}.) Say it once more.
\end{tcolorbox}
\section[\texorpdfstring{\ta{பட்டாம்பூச்சி} (paṭṭāmpūcci)}{பட்டாம்பூச்சி (paṭṭāmpūcci)}]{\ta{பட்டாம்பூச்சி} (paṭṭāmpūcci) — a butterfly}

\label{lesson:TA-C218-pattampucci}



\begin{quote}
Before the new one: say the Tamil for a zebra, then the Tamil for a crocodile.
\end{quote}

\subsection*{You'll want to know: \ta{பட்டாம்பூச்சி}}
\textbf{\ta{பட்டாம்பூச்சி}} — \emph{paṭṭāmpūcci} — \textquotedblleft{}a butterfly\textquotedblright{}.

\begin{cousinweb}
\textbf{\ta{பட்டாம்பூச்சி}} (\emph{paṭṭāmpūcci}) is \emph{paṭṭu}, \textquotedblleft{}silk\textquotedblright{}, from Sanskrit \emph{paṭṭa}, \textquotedblleft{}woven silk\textquotedblright{} (CDIAL 7700), and native \emph{pūcci}, \textquotedblleft{}insect\textquotedblright{} (DEDR 4353): \textquotedblleft{}the silk insect\textquotedblright{}.
\end{cousinweb}

\begin{grammarlens}[title={What you've built}]
One more word for this chapter.
\end{grammarlens}

\subsection*{Your turn}
\begin{itemize}
\raggedright
  \item \emph{Say it:} \emph{paṭṭāmpūcci}
  \item \emph{Say it:} \emph{paṭṭāmpūcci}, once more
  \item \emph{Say it:} say \emph{paṭṭāmpūcci}
  \item \emph{Recall:} say the Tamil for a zebra, then the Tamil for a crocodile, then say \emph{paṭṭāmpūcci} again
\end{itemize}

\begin{tcolorbox}[breakable,colback=teal!4,colframe=teal!35!black,title={Before you move on}]
What does \ta{பட்டாம்பூச்சி} mean? (\textquotedblleft{}A butterfly\textquotedblright{}.) Say it once more.
\end{tcolorbox}
\section[\texorpdfstring{\ta{தேனீ} (tēṉī)}{தேனீ (tēṉī)}]{\ta{தேனீ} (tēṉī) — a bee}

\label{lesson:TA-C218-teni}



\begin{quote}
Before the new one: say the Tamil for a crocodile, then the Tamil for a butterfly.
\end{quote}

\subsection*{You'll want to know: \ta{தேனீ}}
\textbf{\ta{தேனீ}} — \emph{tēṉī} — \textquotedblleft{}a bee\textquotedblright{}.

\begin{cousinweb}
\textbf{\ta{தேனீ}} (\emph{tēṉī}) is native: \emph{tēṉ}, \textquotedblleft{}honey\textquotedblright{} (DEDR 3268), and \emph{ī}, \textquotedblleft{}fly, bee\textquotedblright{} (DEDR 533), which lists \emph{tēṉ-ī}, \textquotedblleft{}honey bee\textquotedblright{}, itself.
\end{cousinweb}

\begin{grammarlens}[title={What you've built}]
One more word for this chapter.
\end{grammarlens}

\subsection*{Your turn}
\begin{itemize}
\raggedright
  \item \emph{Say it:} \emph{tēṉī}
  \item \emph{Say it:} \emph{tēṉī}, once more
  \item \emph{Say it:} say \emph{tēṉī}
  \item \emph{Recall:} say the Tamil for a crocodile, then the Tamil for a butterfly, then say \emph{tēṉī} again
\end{itemize}

\begin{tcolorbox}[breakable,colback=teal!4,colframe=teal!35!black,title={Before you move on}]
What does \ta{தேனீ} mean? (\textquotedblleft{}A bee\textquotedblright{}.) Say it once more.
\end{tcolorbox}
